\section{Boundaries}

RCTX is a bridge, not an ownership transfer. It provides a shared base
dimension that domain packages can reference. It does not own plan assignments,
vote totals, ballot records, audit samples, lineage events, certification
claims, court records, or rendered maps. Those facts remain with RPLAN, RCOUNT,
RHIST, RCERT, RCASE, or RMAP according to the package-family boundary.

The claim boundary is encoded in code as well as prose. \texttt{ClaimBoundary}
requires a package id, a \texttt{proves} list, a \texttt{does\_not\_prove}
list, and optional caveats. \texttt{verify\_package} rejects a package whose
claim boundary names the wrong package or omits either side of the claim. That
shape is modest, but it prevents a shared context package from silently
becoming a legal or cartographic certification.

RMAP is reserved, not implemented. The specs assign RMAP rendered maps,
styling, symbols, labels, layer composition, display projections, tile
products, and map-image hashes. The repository does not currently contain a
standalone RMAP crate, and this paper does not claim one exists. The correct
near-term relationship is that a future RMAP product may reference an RCTX
\texttt{context\_hash}; it should not invent canonical unit identity on its own.

The split decision is also explicit. The plan and boundary spec say not to move
full \texttt{.rctx} read/write out of RPLAN until a second production consumer
needs direct full-context construction or validation. \texttt{rctx-core} exists
because crosswalks, source refs, context unit indexes, graph records, content
hashes, and claim boundaries already have cross-product pressure. That is a
different decision from splitting all RPLAN context I/O.
